The \code{BPatch\_addressSpace} class is a superclass of the \code{BPatch\_process}
and \code{BPatch\_binaryEdit} classes. It contains functionality that is common
between the two subclasses.

\begin{apient}
  BPatch_image *getImage()
\end{apient}

\apidesc{
  Returns a handle to the executable file associated with this \code{BPatch\_process} object.
}

\begin{apient}
  bool getSourceLines(unsigned long addr, std::vector<BPatch_statement> &lines)
\end{apient}

\apidesc{
  Returns the line information associated with the mutatee address, addr. The
  vector lines contain pairs of filenames and line numbers that are associated
  with addr. In many cases only one filename and line number is associated with an
  address, but certain compiler optimizations may lead to multiple filenames and
  lines at an address. This information is only available if the mutatee was compiled
  with debug information.
  
  This function returns true if it was able to find any line information at addr,
  or false otherwise.
}

\begin{apient}
  bool getAddressRanges(const char *fileName, unsigned int lineNo,
                        std::vector<std::pair<unsigned long, unsigned long>> &ranges)
\end{apient}

\apidesc{
  Given a filename and line number, fileName and lineNo, this function this
  function returns the ranges of mutatee addresses that implement the code range
  in the output parameter ranges. In many cases a source code line will only have one
  address range implementing it. However, compiler optimizations may transform this
  into multiple disjoint address ranges. This information is only available if the
  mutatee was compiled with debug information.

  This function returns true if it was able to find any line information, false otherwise.
}

\begin{apient}
  BPatch_variableExpr *malloc(int n, std::string name = std::string(""))

  BPatch_variableExpr *malloc(const BPatch_type &type, std::string name = std::string(""))
\end{apient}

\apidesc{
  These two functions allocate memory. Memory allocation is from a heap. The heap
  is not necessarily the same heap used by the application. The available space in
  the heap may be limited depending on the implementation. The first function
  allocates \code{n} bytes of memory from the heap. The second function allocates
  enough memory to hold an object of the specified type. Using the second version is
  strongly encouraged because it provides additional information to permit better type
  checking of the passed code. If a name is specified, Dyninst will assign \code{var\_name}
  to the variable; otherwise, it will assign an internal name. The returned memory
  is persistent and will not be released until \code{BPatch\_process::free} is
  called or the application terminates.
}

\begin{apient}
  BPatch_variableExpr *createVariable(Dyninst::Address addr, BPatch_type *type,
                                      std::string var_name = std::string(""),
                                      BPatch_module *in_module = NULL)
\end{apient}

\apidesc{
  This method creates a new variable at the given address addr in the module \code{in\_module}.
  If a name is specified, Dyninst will assign \code{var\_name} to the variable; otherwise,
  it will assign an internal name. The type parameter will become the type for the
  new variable.

  When operating in binary rewriting mode, it is an error for the \code{in\_module}
  parameter to be NULL; it is necessary to specify the module in which the
  variable will be created. Dyninst will then write the variable back out in the file
  specified by \code{in\_module}.
}

\begin{apient}
  bool free(BPatch_variableExpr &ptr)
\end{apient}

\apidesc{
  Free the memory in the passed variable ptr. The programmer is responsible for verifying
  that all code that could reference this memory will not execute again (either by
  removing all snippets that refer to it, or by analysis of the program). Return true
  if the free succeeded.
}

\begin{apient}
  bool getRegisters(std::vector<BPatch_register> &regs)
\end{apient}

\apidesc{
  Returns the registers available to snippet code.
}

\begin{apient}
  BPatchSnippetHandle *insertSnippet(const BPatch_snippet &expr,  BPatch_point &point,
                                     BPatch_callWhen when,
                                     BPatch_snippetOrder order = BPatch_firstSnippet)

  BPatchSnippetHandle *insertSnippet(const BPatch_snippet &expr,
                                     const std::vector<BPatch_point *> &points,
                                     BPatch_callWhen when, 
                                     BPatch_snippetOrder order = BPatch_firstSnippet)
\end{apient}

\apidesc{
  Insert a snippet of code at the specified point. If a list of points is supplied,
  insert the code snippet at each point in the list. The optional when argument specifies
  when the snippet is to be called; a value of \code{BPatch\_callBefore} indicates
  that the snippet should be inserted just before the specified point or points in
  the code, and a value of \code{BPatch\_callAfter} indicates that it should be inserted
  just after them.

  The order argument specifies where the snippet is to be inserted relative to any
  other snippets previously inserted at the same point. The values \code{BPatch\_firstSnippet}
  and \code{BPatch\_lastSnippet} indicate that the snippet should be inserted
  before or after all snippets, respectively.

  It is illegal to use \code{BPatch\_callAfter} with a \code{BPatch\_entry} point.
  Use \code{BPatch\_callBefore} when instrumenting entry points, which inserts
  instrumentation before the first instruction in a subroutine. Likewise, it is
  illegal to use \code{BPatch\_callBefore} with a \code{BPatch\_exit} point. Use
  \code{BPatch\_callAfter} with exit points. \code{BPatch\_callAfter} inserts
  instrumentation at the last instruction in the subroutine. insert­Snippet will return
  NULL when used with an illegal pair of points.
}

\begin{apient}
  bool deleteSnippet(BPatchSnippetHandle *handle)
\end{apient}

\apidesc{
  Remove the snippet associated with the passed handle. If the handle is not defined
  for the process, then deleteSnippet will return false.
}

\begin{apient}
  void beginInsertionSet()
\end{apient}

\apidesc{
  Normally, a call to insertSnippet immediately injects instrumentation into the mutatee.
  However, users may wish to insert a set of snippets as a single batch operation.
  This provides two benefits: First, Dyninst may insert instrumentation in a more
  efficient manner. Second, multiple snippets may be inserted at multiple points
  as a single operation, with either all snippets being inserted successfully or none.
  This batch insertion mode is begun with a call to beginInsertionSet; after this call,
  no snippets are actually inserted until a corresponding call to
  finalizeInsertionSet. Dyninst accumulates all calls to insertSnippet during batch
  mode internally, and the returned BPatchSnippetHandles are filled in when
  finalizeInsertionSet is called.

  Insertion sets are unnecessary when doing static binary instrumentation. Dyninst
  uses an implicit insertion set around all instrumentation to a static binary.
}

\begin{apient}
  bool finalizeInsertionSet(bool atomic)
\end{apient}

\apidesc{
  Inserts all snippets accumulated since a call to beginInsertionSet. If the atomic
  parameter is true, then a failure to insert any snippet results in all snippets
  being removed; effectively, the insertion is all-or-nothing. If the atomic parameter
  is false, then snippets are inserted individually. This function also fills in the
  BPatchSnippetHandle structures returned by the insertSnippet calls comprising this
  insertion set. It returns true on success and false if there was an error
  inserting any snippets.

  Insertion sets are unnecessary when doing static binary instrumentation. Dyninst
  uses an implicit insertion set around all instrumentation to a static binary.
}

\begin{apient}
  bool removeFunctionCall(BPatch_point &point)
\end{apient}

\apidesc{
  Disable the mutatee function call at the specified location. The point specified
  must be a valid call point in the image of the mutatee. The purpose of this
  routine is to permit tools to alter the semantics of a program by eliminating
  procedure calls. The mechanism to achieve the removal is platform dependent, but
  might include branching over the call or replacing it with NOPs. This function only
  removes a function call; any parameters to the function will still be evaluated.
}

\begin{apient}
  bool replaceFunction(BPatch_function &old, BPatch_function &new)

  bool revertReplaceFunction(BPatch_function &old)
\end{apient}

\apidesc{
  Replace all calls to user function old with calls to new. This is done by inserting
  instrumentation (specifically a \code{BPatch\_funcJumpExpr}) into the beginning
  of function old such that a non-returning jump is made to function new. Returns true
  upon success, false otherwise.
}

\begin{apient}
  bool replaceFunctionCall(BPatch_point &point, BPatch_function &newFunc)
\end{apient}

\apidesc{
  Change the function call at the specified point to the function indicated by
  newFunc. The purpose of this routine is to permit runtime steering tools to change
  the behavior of programs by replacing a call to one procedure by a call to another.
  Point must be a function call point. If the change was successful, the return
  value is true, otherwise false will be returned.

  \it{\em{WARNING}}: Care must be used when replacing functions. In particular if the
  compiler has performed inter-procedural register allocation between the original
  caller\/callee pair, the replacement may not be safe since the replaced function may
  clobber registers the compiler thought the callee left untouched. Also the signatures
  of the both the function being replaced and the new function must be compatible.
}

\begin{apient}
  bool wrapFunction(BPatch_function *old, BPatch_function *new, Dyninst::SymtabAPI::Symbol *sym)

  bool revertWrapFunction(BPatch_function *old)
\end{apient}

\apidesc{
  Replaces all calls to function old with calls to function new. Unlike replaceFunction
  above, the old function can still be reached via the name specified by the
  provided symbol sym. Function wrapping allows existing code to be extended by new
  code. Consider the following code that implements a fast memory allocator for a
  particular size of memory allocation, but falls back to the original memory allocator
  (referenced by origMalloc) for all others.
}

  \begin{lstlisting}[language=C++]
        void *origMalloc(unsigned long size);
    
        void *fastMalloc(unsigned long size) {
          if (size == 1024) {
            unsigned long ret = fastPool;
            fastPool += 1024;
            return ret;
          }
          else {
            return origMalloc(size);
          }
        }
  \end{lstlisting}

\apidesc{
  The symbol sym is provided by the user and must exist in the program; the easiest
  way to ensure it is created is to use an undefined function as shown above with the
  definition of origMalloc.

  The following code wraps malloc with fastMalloc, while allowing functions to
  still access the original malloc function by calling origMalloc. It makes use of
  the new convert interface described in Section 5.
}

  \begin{lstlisting}[language=C++]
        using namespace Dyninst;
        using namespace SymtabAPI;
    
        BPatch_function *malloc = appImage->findFunction(...);
        BPatch_function *fastMalloc = appImage->findFunction(...);
        Symtab *symtab = SymtabAPI::convert(fastMalloc->getModule());
        std::vector<Symbol *> syms;
        symtab->findSymbol(syms, "origMalloc",
          Symbol::ST_UNKNOWN, // Don't specify type
          mangledName,        // Look for raw symbol name
          false,              // Not regular expression
          false,              // Don't check case
          true                // Include undefined symbols
        );
        app->wrapFunction(malloc, fastMalloc, syms[0]);
  \end{lstlisting}

\apidesc{
  For a full, executable example, see Appendix A Complete Examples.
}

\begin{apient}
  bool replaceCode(BPatch_point *point, BPatch_snippet *snippet)
\end{apient}

\apidesc{
  This function has been removed; users interested in replacing code should instead use
  the PatchAPI code modification interface described in the PatchAPI manual. For
  information on accessing PatchAPI abstractions from DyninstAPI abstractions, see Section 5.
}

\begin{apient}
  BPatch_module * loadLibrary(const char *libname, bool reload=false)
\end{apient}

\apidesc{
  For dynamic rewriting, this function loads a dynamically linked library into the
  process\'s address space. For static rewriting, this function adds a library as a
  library dependency in the rewritten file. In both cases Dyninst creates a new
  \code{BPatch\_module} to represent this library.

  The libname parameter identifies the file name of the library to be loaded, in the
  standard way that dynamically linked libraries are specified on the operating system
  on which the API is running. This function returns a handle to the loaded
  library. The reload parameter is ignored and only remains for backwards compatibility.
}

\begin{apient}
  bool isStaticExecutable()
\end{apient}

\apidesc{
  This function returns true if the original file opened with this
  \code{BPatch\_addressSpace} is a statically linked executable, or false otherwise.
}

\begin{apient}
  processType getType()
\end{apient}

\apidesc{
  This function returns a processType that reflects whether this address space is
  a \code{BPatch\_process} or a \code{BPatch\_binaryEdit}.
}
